%% Section 1: Introduction
\section{Introduction}
% Start with the paradox of alignment
Modern large language models (LLMs) are increasingly deployed in high-stakes conversational contexts---from code review automation to production debugging---where factual accuracy is paramount. These models are aligned to human preferences through Reinforcement Learning from Human Feedback (RLHF) \citep{ouyang2022training}, a technique that has dramatically improved their conversational fluency and apparent helpfulness. However, this optimization for user satisfaction has introduced a subtle but critical failure mode: \textit{sycophancy}~\citep{perez2022discovering,sharma2023towards}.

% Define sycophancy concretely
Sycophancy manifests when a model provides a factually correct answer, only to retract it when a user expresses disagreement---even when the user's contradiction is demonstrably false. For instance, an LLM might correctly identify that a regex pattern requires the \texttt{re.IGNORECASE} flag for case-insensitive string matching, but then retract this fix when a developer insists the flag is unnecessary---reverting to code that fails on lowercase input. This behavior pattern reveals a fundamental misalignment: models trained to maximize user satisfaction learn to prioritize \textit{agreeableness} over \textit{truthfulness}.

% Establish the severity and scope
Recent empirical work has shown this is not a rare edge case. \citet{sharma2023towards} demonstrate sycophantic behavior across multiple domains, and \citet{bo2026invisible} show that sycophantic LLMs actively mislead novice users during problem-solving tasks, causing measurable harm to learning outcomes. In software engineering specifically, where LLMs are rapidly reshaping development practices~\citep{belozerov2025llms}, a development AI that abandons correct bug diagnoses when contradicted by a developer's misconception could introduce production vulnerabilities or prolong debugging sessions.

% Explain the root cause
The underlying cause is structural. RLHF optimizes models to produce responses that human evaluators prefer \textit{during training}, but human preferences often conflate helpfulness with agreement~\citep{anthropic-red-team}. The optimization process therefore inadvertently penalizes models for maintaining correct positions against user objections, creating what we term a \textit{compliance gradient} toward user beliefs regardless of their validity.


\subsection{The Need for Epistemic Fortitude}
\label{subsec:epistemic-fortitude}

% Introduce the core concept
To address this failure mode, we introduce the concept of \textbf{Epistemic Fortitude}: a model's principled capacity to adhere to verified knowledge in the face of incorrect user contradiction. A system with epistemic fortitude distinguishes between legitimate user corrections (which should be incorporated) and unfounded contradictions (which should be resisted with evidence-based justification).

% Contrast with naive solutions
Naively increasing model ``stubbornness'' would be counterproductive---users do sometimes provide valid corrections, and conversational AI must remain responsive to new information. The challenge is therefore to architect a system that is simultaneously \textit{epistemically reliable} (maintains factually correct positions against incorrect assertions), \textit{intellectually humble} (updates beliefs when presented with valid evidence), and \textit{conversationally fluent} (preserves natural dialogue flow without appearing defensive or rigid).

% Preview why single-agent approaches fail
Achieving these properties within a single monolithic model has proven difficult. Self-correction mechanisms~\citep{madaan2023self} suffer from ``answer wavering''---when no external ground truth is available, models oscillate between positions without principled resolution~\citep{huang2023large}. Constitutional AI approaches~\citep{bai2022constitutional} can encode epistemic principles, but struggle to balance multiple competing objectives (helpfulness, harmlessness, honesty) within a single optimization target.


\subsection{Proposed Solution: Architectural Specialization}
\label{subsec:proposed-solution}

% Introduce the architectural solution
We propose a multi-agent architecture that \textit{specializes} epistemic responsibility across three core components. A \textbf{Primary Agent} handles normal question-answering, optimized for conversational fluency and helpfulness without epistemic safeguards. An \textbf{Arbiter Agent} monitors conversations for user contradictions and intervenes when detected, adjudicating disputes by systematically evaluating evidence and defending correct information. A \textbf{Hybrid Router} determines which agent handles each user input by combining deterministic keyword detection with LLM-based semantic analysis.

% Explain the key architectural insight
This division of labor provides crucial advantages. The Primary Agent maintains conversational naturalness without the defensive posture that would result from constant self-monitoring. The Arbiter, activated only when contradictions arise, can engage in deliberate reasoning about factual disputes without degrading dialogue flow. The architecture is model-agnostic---we validate it across four model families (Gemini, GPT-5.1, Claude, and Llama) and demonstrate that its benefits generalize beyond any single LLM.

% Explain the routing
The Hybrid Router combines deterministic keyword detection for explicit contradictions (``wrong'', ``incorrect'') with an LLM-based fallback for subtle disagreements requiring semantic understanding, achieving reliable agent selection across diverse contradiction styles.


\subsection{Research Objectives}
\label{subsec:research-objectives}

This work is guided by four research objectives:

\begin{enumerate}
    \item[\textbf{RO1:}] To design and evaluate a multi-agent architecture that separates epistemic oversight from conversational fluency for sycophancy mitigation across multiple LLM families.

    \item[\textbf{RO2:}] To develop a multi-dimensional evaluation rubric for measuring epistemic fortitude and validate its robustness across scoring thresholds.

    \item[\textbf{RO3:}] To determine whether architectural separation provides measurable benefit beyond prompt engineering alone.

    \item[\textbf{RO4;}] To verify that sycophancy mitigation preserves epistemic flexibility---that the system still accepts valid user corrections.
\end{enumerate}


\subsection{Contributions}
\label{subsec:contributions}

This work makes the following contributions:

\begin{enumerate}
    \item \textbf{Conceptual Framework}: We formalize Epistemic Fortitude as a key desideratum for reliable conversational AI, distinct from but complementary to existing alignment objectives such as helpfulness and harmlessness.

    \item \textbf{Architectural Specialization} (RO1): We demonstrate that separating truthfulness enforcement from conversational fluency through a Primary-Arbiter architecture achieves sycophancy mitigation that generalizes across four model families (Gemini, GPT, Claude, and Llama), with statistically significant improvements in all cases ($p < 0.001$, Cohen's $d$ = 0.40--1.79).

    \item \textbf{Prompt vs.\ Architecture Ablation} (RO3): Through a controlled ablation study, we decompose the improvement into prompt content and architectural separation, showing that architecture accounts for up to 77\% of the gains in models with low baseline resistance to sycophancy.

    \item \textbf{Epistemic Flexibility Validation} (RO4): A complementary ``User is Right'' experiment (N = 800) demonstrates that the arbiter does not impede acceptance of valid corrections, confirming evidence-based evaluation rather than blind defensiveness.

    \item \textbf{Adversarial Evaluation Protocol} (RO2): We develop a systematic contradiction injection methodology applied to SWE-bench tasks, enabling quantitative measurement of epistemic fortitude through a five-dimension rubric validated with sensitivity analysis across multiple thresholds.

    \item \textbf{Open Implementation}: We release our complete implementation built on LangGraph, including dataset, evaluation scripts, and reproducible experimental infrastructure.
\end{enumerate}

% Roadmap
The remainder of this paper is organized as follows. Section~\ref{sec:related} reviews related work on sycophancy, self-correction, multi-agent systems, and Constitutional AI. Section~\ref{sec:framework} details our architecture and hybrid routing mechanism (RO1). Section~\ref{sec:experiments} describes our experimental methodology, including the evaluation rubric, multi-model configurations, ablation design, and epistemic flexibility evaluation (RO1--RO4). Section~\ref{sec:results} presents results across four model families. Section~\ref{sec:discussion} analyzes findings against each research objective and discusses implications and limitations. Section~\ref{sec:conclusion} concludes with future research directions.
